\chapter{Comprehensive Conclusions and Subsequent Validation}
\section{Established Control and Validation Pipeline}
\subsection{Trajectory and Feedback System}
The target, state, and clock interfaces between the ROS2 control node and the Isaac physics simulation have been established. To address initial tracking lag, adjusting the reference velocity and drive gain under independent motion amplitudes reduced the average motion error from 0.864 degrees to 0.082 degrees. This result is limited to the same Panda model and scenario and cannot be directly applied to another robotic arm.

The state feedback portion connects the virtual IMU, publicly available real-world data, and Jetson hardware-based acquisition. The most explicit engineering conclusion is that information freshness differs from publication frequency: physical sensors output at approximately 25 Hz, while the majority of published messages at 200 Hz reuse measurements; five-point averaging introduces approximately 80 ms group delay. The 95th percentile of message transmission latency being under 1 ms does not imply that the entire measurement-to-execution pipeline is under 1 ms.
\subsection{Contact Control}
Normal admittance can reduce sustained pressing force, but the error integral of direct lift force during flight will accelerate downward pressure, thereby increasing the initial impact. Contact gating, approach deceleration, target force ramp-up, continuous inverse dynamics, and gravity compensation address different stages separately and should not be reduced to the effect of a single parameter.

The later-stage selection criteria, under a 1 ms time step, show an original peak force of 11.11 N and an average force of approximately 5.007 N in the final phase, satisfying the four predefined project checks; this is not an industrial certification standard. Reducing the step size by half under rigid contact increases the peak force to 20.43 N, and certain height and friction conditions also reveal sensitivity. Soft material improves the numerical stability of this set of values, but changing the environment cannot be considered as the rigid controller itself having resolved all issues.
\subsection{From Simplified Balance to Realistic Multi-Joint Walking}
The linear inverted pendulum provides constrained support point and footstrike analysis, while the finite-foot model supplements realistic contact. Subsequently, the six-actuator planar biped connects leg dynamics, whole-body inverse dynamics, and real foot-ground contact, advancing approximately 0.937 meters under the eight-step standard condition; later, both groups of twenty steps advance approximately 2.323 meters.

Under explicit zero-moment-point task conditions at 100 Hz with no disturbances, the actual pressure center error was reduced from 19.41 mm to 6.06 mm, but the center-of-mass error increased, and backward recovery became slower. Among 94 unique conditions, 79 passed, but failures due to long delays still occurred frequently. This indicates that even after unifying the physical quantity interface, trade-offs remain between multi-task weighting, contact switching, and state prediction.
\subsection{Whole-body Control}
The G1 model retains a free-floating base, realistic collisions, and 29 constrained motors, enabling single-foot support on the left and right sides and three-dimensional slow alternating stepping. In four steps, it advances approximately 0.210 meters in 29 seconds; in six steps, it advances approximately 0.330 meters in 43 seconds. These results exceed planar standing but remain within the realm of quasistatic walking, without validation of dynamic high-speed motion.

Out of 52 test conditions, 51 passed, while the remaining 8 under strong disturbances only passed 2. Solver residual thresholds, self-contact, and support phase timing are all actual limitations. Clarabel exhibits relatively short tail latency in call time across the tested conditions, but does not include all sensing, communication, and scheduling overheads, and thus cannot be termed a hard real-time controller based on this.
\section{Engineering Experience Gained Across Stages}
\begin{enumerate}
\item First, verify the units, symbols, and execution channels. The ground support point is in meters, and the contact wrench and motor torque are not the same variable; they can only be converted into execution commands after full dynamic constraints are satisfied.
\item Proceed layer by layer under the assumption of increasing realism. Disabling gravity, simplifying collisions, and restricting planar degrees of freedom help isolate issues, but revalidation is required once these conditions are restored.
\item Both average metrics and tail risks are preserved. Original force peaks, tail-end solution times, early terminations, and contact confirmations cannot be masked by averages.
\item The controller can only use information available at the current time. Delayed queues, frame reuse, drift, and inputs to prediction models must be explicitly defined and must not covertly read the true current state.
\item Separate numerical verification from task acceptance. Consistent contact force recalculation and small dynamics residuals indicate self-consistency; whether the walk is completed, whether weight-bearing is restored, and whether timing constraints are satisfied must each be judged independently.
\end{enumerate}
\section{Deployment Conditions Still to be Completed}
\subsection{State Estimation and Contact Events}
Integrate IMU, joint encoders, and foot contact into a single state estimator, explicitly accounting for drift, bias, and packet loss; adjust the nominal timing based on actual touchdown events. Current measurement testing and control simulation are conducted separately; the real-sensor closed loop has not yet been completed.
\subsection{Actuators and Additional Contacts}
Limited torque does not equal the real motor model. In the next phase, bandwidth, communication, rate, and thermal constraints should be incorporated, along with control constraints for self-collision, premature touchdown, and terrain variation. The observed wrist-hip and shoulder-trunk contacts under strong disturbances provide specific points for investigation, but they cannot yet be identified as the sole cause of all failures.
\subsection{Independent Verification of Contact Control}
Prioritize validation of unknown surface heights, slopes, tangential motion, exit and recovery after loss of contact, and verify force peaks using calibrated finite stiffness and measured bandwidth. Fixed-plane single normal results do not automatically generalize to six-dimensional compliant operation.
\subsection{Preparation for Hardware Testing}
If suitable equipment is obtained, first confirm the robot's torque interface, built-in gravity compensation, state timestamp, and stopping mechanism, then select low-risk free-space trajectories and limited contact conditions. The test materials should include interface specifications, parameter limits, data recording, and clearly defined acceptance criteria. Currently, no proxy test results have been obtained, and the scheme preparation has not been counted as a successful real hardware pass.
\section{How to Use the Report and Code}
The source code and original reports for each phase are indexed in Appendix A; model sources, configurations, and dependencies are retained in their original positions. Appendix B lists copied images and English-named videos, Appendix C explains terminology, and Appendix D preserves historical run examples. The source code directory does not constitute an independent software distribution package containing all assets and environments; to reproduce, first read the corresponding README, and establish a new output directory for each new run.

\section{Official Walking Baseline Supplementary Conclusions}
The dynamic walking of the official reinforcement learning policy has been reproduced; under a 0.5 m/s condition, the robot advanced 9.23 meters and drifted 1.41 meters over 20 seconds, with a heading change of approximately 15.6 degrees. After zeroing the command, average net displacement slowed but the robot continued stepping, failing to meet the preset instantaneous stop condition. The self-developed whole-body controller maintained quasi-static prototype positioning; after this phase, further development of natural gaits will not be pursued, with mode switching, heading maintenance, state estimation, and hardware validation retained as clear boundaries.
